\documentclass[10pt,a4paper]{amsart}
\usepackage{amsmath,amssymb,amsthm,graphicx}
% For commutative diagrams we use Xy-pic
\usepackage[all]{xy}

% We use 2cell for 2-commutative diagrams.
\xyoption{2cell}
% Theorem environments.
%
\theoremstyle{plain}
\newtheorem{theorem}[subsection]{Theorem}
\newtheorem{proposition}[subsection]{Proposition}
\newtheorem{lemma}[subsection]{Lemma}

\theoremstyle{definition}
\newtheorem{definition}[subsection]{Definition}
\newtheorem{example}[subsection]{Example}
\newtheorem{exercise}[subsection]{Exercise}
\newtheorem{situation}[subsection]{Situation}

\theoremstyle{remark}
\newtheorem{remark}[subsection]{Remark}
\newtheorem{remarks}[subsection]{Remarks}

\numberwithin{equation}{subsection}

% Macros
%
\def\lim{\mathop{\mathrm{lim}}\nolimits}
\def\colim{\mathop{\mathrm{colim}}\nolimits}
\def\Spec{\mathop{\mathrm{Spec}}}
\def\Hom{\mathop{\mathrm{Hom}}\nolimits}
\def\Ext{\mathop{\mathrm{Ext}}\nolimits}
\def\SheafHom{\mathop{\mathcal{H}\!\mathit{om}}\nolimits}
\def\SheafExt{\mathop{\mathcal{E}\!\mathit{xt}}\nolimits}
\def\Sch{\mathit{Sch}}
\def\Mor{\mathop{\mathrm{Mor}}\nolimits}
\def\Ob{\mathop{\mathrm{Ob}}\nolimits}
\def\Sh{\mathop{\mathit{Sh}}\nolimits}
\def\NL{\mathop{N\!L}\nolimits}
\def\CH{\mathop{\mathrm{CH}}\nolimits}
\def\proetale{{pro\text{-}\acute{e}tale}}
\def\etale{{\acute{e}tale}}
\def\QCoh{\mathit{QCoh}}
\def\Ker{\mathop{\mathrm{Ker}}}
\def\Im{\mathop{\mathrm{Im}}}
\def\Coker{\mathop{\mathrm{Coker}}}
\def\Coim{\mathop{\mathrm{Coim}}}

% Boxtimes
%
\DeclareMathSymbol{\boxtimes}{\mathbin}{AMSa}{"02}

%
% Macros for moduli stacks/spaces
%
\def\QCohstack{\mathcal{QC}\!\mathit{oh}}
\def\Cohstack{\mathcal{C}\!\mathit{oh}}
\def\Spacesstack{\mathcal{S}\!\mathit{paces}}
\def\Quotfunctor{\mathrm{Quot}}
\def\Hilbfunctor{\mathrm{Hilb}}
\def\Curvesstack{\mathcal{C}\!\mathit{urves}}
\def\Polarizedstack{\mathcal{P}\!\mathit{olarized}}
\def\Complexesstack{\mathcal{C}\!\mathit{omplexes}}
% \Pic is the operator that assigns to X its picard group, usage \Pic(X)
% \Picardstack_{X/B} denotes the Picard stack of X over B
% \Picardfunctor_{X/B} denotes the Picard functor of X over B
\def\Pic{\mathop{\mathrm{Pic}}\nolimits}
\def\Picardstack{\mathcal{P}\!\mathit{ic}}
\def\Picardfunctor{\mathrm{Pic}}
\def\Deformationcategory{\mathcal{D}\!\mathit{ef}}
\usepackage{cleveref}
\crefname{originallemma}{original lemma}{original lemmas}
\Crefname{originallemma}{Original lemma}{Original lemmas}
\makeatletter
\newlabel{modules-lemma-support-section-closed}{{01AU}{source}{}{sourcepointer-1}{}}
\newlabel{modules-lemma-support-section-closed@cref}{{[originallemma][0][]01AU}{[1][0][]source}{}{}{}}
\makeatother
\title{The locus where a local section of an etale sheaf vanishes is the largest Zariski open subset and is detected by geometric points}
\author{The Stacks Project Authors}
\date{stacks-project:etale-cohomology.tex:local-snapshot-20261011; GNU FDL 1.2 or later}
\begin{document}
\maketitle
\noindent\textit{Editorial note.} All delivered mathematical text is the author's own, copied from the pinned Stacks Project chapter file etale-cohomology.tex; no statement or proof step was rewritten, reordered or removed.\par\medskip
\noindent\textit{Editorial note.} The proof is the author's complete argument for the three numbered parts, each written out inside the counted span.\par\medskip
\noindent\textit{Editorial note.} The label lemma-zero-over-image occurs in no other staged Stacks chapter, so no cross-chapter restatement had to be weighed for this card.\par\medskip
\noindent\textit{Editorial note.} The single labelled reference occurring inside the counted spans points to another chapter of the Project; it is carried in assembly.external\_references as a non-load-bearing navigation pointer recorded at the invoking line.\par\medskip
\noindent\textit{Editorial note.} The preamble is the chapter's own native preamble excerpt already used for this source: the Xy-pic setup (native preamble lines 14-18) and the native theorem declarations with the mathematical macros (native preamble lines 154-217).\par\medskip
\noindent\textit{Allowed standard prerequisites (editorial).}\begin{itemize}
\item Restriction of an etale sheaf to the Zariski site of a scheme. Same chapter, the fact that a sheaf for the etale topology restricts to a sheaf for the Zariski topology on U. Supplies the Zariski sheaf whose vanishing locus W is produced in part (1).
\item Sheaf condition and etale coverings V -> phi(V). Same chapter, the sheaf condition for F on U\_etale applied to the etale covering \{V -> phi(V)\}. Supplies the step that a section vanishing on V vanishes on the image phi(V), which proves part (2).
\item Stalk of an etale sheaf at a geometric point and etale neighbourhoods. Same chapter, the description of the stalk F\_sbar as classes of sections over etale neighbourhoods, with vanishing meaning vanishing on some etale neighbourhood. Supplies the equivalence in part (3) between vanishing at a geometric point and containment of that point in W.
\end{itemize}\medskip
\ifdefined\hypertarget\hypertarget{sourcepointer-1}{}\fi
\noindent\textit{Source reference.} Original lemma 01AU: sources/stacks-etale-cohomology/pinned-original-source.tex, source lines 3969--3969. Cross-chapter navigation pointer invoked at proof lines 3968-3969, where the author uses that the zero locus of a section of a sheaf on a scheme is Zariski open. The target chapter is not staged in this card's source tree, so the pointer records the invoking line; it only resolves the reference for the reader, and no bibliography citation is involved.\par\medskip
\begin{lemma}
\label{lemma-zero-over-image}
Let $S$ be a scheme.
Let $\mathcal{F}$ be an abelian sheaf on $S_\etale$.
Let $\sigma \in \mathcal{F}(U)$ be a local section.
There exists an open subset $W \subset U$ such that
\begin{enumerate}
\item $W \subset U$ is the largest Zariski open subset of $U$ such
that $\sigma|_W = 0$,
\item for every $\varphi : V \to U$ in $S_\etale$ we have
$$
\sigma|_V = 0 \Leftrightarrow \varphi(V) \subset W,
$$
\item for every geometric point $\overline{u}$ of $U$ we have
$$
(U, \overline{u}, \sigma) = 0\text{ in }\mathcal{F}_{\overline{s}}
\Leftrightarrow
\overline{u} \in W
$$
where $\overline{s} = (U \to S) \circ \overline{u}$.
\end{enumerate}
\end{lemma}

\begin{proof}
Since $\mathcal{F}$ is a sheaf in the \'etale topology the restriction of
$\mathcal{F}$ to $U_{Zar}$ is a sheaf on $U$ in the Zariski topology.
Hence there exists a Zariski open $W$ having property (1), see
Modules, Lemma \ref{modules-lemma-support-section-closed}. Let
$\varphi : V \to U$ be an arrow of $S_\etale$. Note that
$\varphi(V) \subset U$ is an open subset and that
$\{V \to \varphi(V)\}$ is an \'etale covering. Hence if
$\sigma|_V = 0$, then by the sheaf condition for $\mathcal{F}$ we
see that $\sigma|_{\varphi(V)} = 0$. This proves (2).
To prove (3) we have to show that if $(U, \overline{u}, \sigma)$
defines the zero element of $\mathcal{F}_{\overline{s}}$, then
$\overline{u} \in W$. This is true because the assumption means
there exists a morphism of \'etale neighbourhoods
$(V, \overline{v}) \to (U, \overline{u})$ such that
$\sigma|_V = 0$. Hence by (2) we see that $V \to U$ maps into $W$, and
hence $\overline{u} \in W$.
\end{proof}

\end{document}
